\documentclass[../../main.tex]{subfiles}
\begin{document}

\begin{center}
  {\fontsize{48pt}{58pt}\selectfont \textbf{1974}\normalsize\textbf{年度}}
\end{center}

\begin{table}[htb]
  \centering
  \caption{出題テーマと難易度}
  \begin{tabular}{|c|p{9cm}|c|}
    \hline
    \textbf{問題番号} & \textbf{テーマ} & \textbf{難易度} \\
    \hline
    第1問 & 二次関数の値域 & \\
    \hline
    第2問 & 複素数と方程式の根 & \\
    \hline
    第3問 & 複素数と写像の値域 & \\
    \hline
    第4問 & 楕円と回転体の体積 & \\
    \hline
    第5問 & 指数関数の最大値 & \\
    \hline
    第6問 & 定積分と絶対値関数の最小値 & \\
    \hline
  \end{tabular}
\end{table}

\begin{table}[htb]
  \centering
  \caption{解答}
  \renewcommand{\arraystretch}{1.4}
  \begin{tabular}{|c|c|p{9cm}|}
    \hline
    \textbf{問題番号} & \textbf{小問} & \textbf{答え} \\
    \hline
    第1問 & - & $(a,b)=(1,4)$ \\
    \hline
    第2問 & (1) & 証明問題 \\
    \cline{2-3}
    & (2) & $f(x)=x^3-x,\ x^3+x^2+x,\ x^3-1$ \\
    \hline
    第3問 & (1) & $x=-\dfrac{b}{2a}$ または $y=0$ \\
    \cline{2-3}
    & (2) & 証明問題 \\
    \hline
    第4問 & - & $V=\sqrt{3}\pi a^2b$ \\
    \hline
    第5問 & - & $a=e\pm2$ \\
    \hline
    第6問 & - & $a=\dfrac{\sqrt{3}}{3}$のとき最小，$F\left(\dfrac{\sqrt{3}}{3}\right)=\sqrt{3}-1$ \\
    \hline
  \end{tabular}
\end{table}

\paragraph{試験形態}
6問．試験時間は要確認．

\end{document}
